% Member 4 - final report, III.C Structure-Preserving Style Augmentation
\subsection{Structure-Preserving Style Augmentation}
\label{sec:structure_method}

Our method changes one thing relative to AdaIN: the constant strength in
Equation~\eqref{eq:blend} becomes a field that depends on the local
structure of the content image. This section describes how that field is
computed.

\paragraph{The structure map.}
Let $M(x) \in [0,1]^{H \times W}$ be a soft structural response of the
content image $x$. It is computed with classical image processing, with no
learned segmentation model and no target-domain information of any kind.
The content image is converted to greyscale and blurred with a Gaussian of
$\sigma = 1$; the horizontal and vertical Sobel responses are combined into
a gradient magnitude $G$\cite{opencv_sobel}. To make the response comparable across images, it
is normalised by the 95th percentile of $G$ over the image and clipped to
$[0,1]$; when that percentile is zero the maximum is used instead, and a
constant image yields an all-zero map. The normalised response is then
dilated with a $3\times3$ max-pooling and smoothed with a Gaussian of
$\sigma = 2$, which widens thin contours and makes the resulting field
smooth enough to be used as a blending weight. The result is $M \in [0,1]$
with $M \approx 1$ at strong contours and $M \approx 0$ in flat regions.

\paragraph{The strength field.}
The field itself is a linear function of $M$:
\begin{equation}
\alpha(x)(h,w) = \alpha_0 \bigl(1 - \beta\, M(x)(h,w)\bigr),
\label{eq:alphamap}
\end{equation}
with $\alpha_0 \in [0,1]$ the maximum strength and $\beta \in [0,1]$ the
modulation depth. Where the structure is strong the strength falls towards
$\alpha_0(1-\beta)$; where it is weak the strength is $\alpha_0$. Two
boundary cases make the design easy to reason about. At $\beta = 0$ the
field is constant and Equation~\eqref{eq:blend} reduces exactly to AdaIN at
strength $\alpha_0$: the baseline is a special case of the method, not a
separate implementation. At $\alpha_0 = 0$ the output is the
encoder--decoder reconstruction.

\paragraph{Why the field is not a segmentation.}
$M$ is an edge response, not a foreground mask. Backgrounds have edges
too, and they are protected as well; conversely, a region with weak texture
inside the object receives the full strength. This is a deliberate choice.
A learned segmentation would be a stronger prior, but it would introduce a
model whose own errors would have to be controlled, and it would risk
leaking information about the classes. The edge response has neither
problem, and Section~\ref{sec:strength_relation} shows what it costs: the
spatial structure of the field turns out not to matter for the accuracy.

\paragraph{Resolution and its consequences.}
The field is computed at the resolution of the content image and then
area-averaged down to the resolution of the $\texttt{relu4\_1}$ feature map
on which the blend is performed, so that one weight corresponds to one
feature location. The average is an area average, not a subsampling, so the
mean of the field over the feature grid is preserved. This is what makes the
comparison of mean strengths exact and not merely approximate, and it is
why $\bar{\alpha}$ of Equation~\eqref{eq:meanstrength} can be computed and
reported for every configuration.

\paragraph{Implementation.}
The generator reuses the pretrained AdaIN encoder--decoder without any
training: the same network, the same weights, the same preprocessing (the
content image is resized so that its short side is 512, preserving the
aspect ratio). Only the blending rule differs. All the content--style
pairs, the image sizes, the file format and the JPEG quality are identical
across every configuration, so that the only variable between an
experimental group and its control is the rule that computes $\alpha$.
